\documentclass[a4paper,11pt]{article}
\usepackage[utf8]{inputenc}
\usepackage[T1]{fontenc}
\usepackage[french]{babel}
\usepackage[left=2cm,right=2cm,top=2cm,bottom=2cm]{geometry}
\usepackage{xcolor}
\usepackage{hyperref}
\usepackage{fontawesome5}
\usepackage{parskip}

% --- Couleurs ---
\definecolor{primary}{RGB}{35, 55, 123}

\begin{document}

% --- EN-TÊTE ---
\begin{center}
    {\Large \textbf{Loïc Aron MBASSI EWOLO}}\\[4pt]
    \small
    \faMapMarker\ Cergy, France \quad
    \faPhone\ (+33) 7 59 52 33 98 \quad
    \faEnvelope\ \href{mailto:loicaron.mbassiewolo@ensea.fr}{loicaron.mbassiewolo@ensea.fr}
\end{center}
\vspace{0.5cm}

% --- DESTINATAIRE ---
\hfill
\begin{minipage}[t]{0.45\textwidth}
    \textbf{CEA Gramat}\\
    Service Recrutement\\
    Gramat (46)
\end{minipage}

\vspace{1cm}

\textbf{Objet :} Candidature Stage Développement d'outils d'IA pour l'analyse de défauts de faradisation

\vspace{0.5cm}

Madame, Monsieur,

La détection automatique d'anomalies et le traitement du signal sont des sujets sur lesquels j'ai travaillé concrètement durant mon parcours. Votre offre de stage sur le développement d'outils d'IA pour analyser les défauts de cages de Faraday correspond à mes compétences en apprentissage automatique et en Python. Actuellement en double diplôme à l'\textbf{ENSEA} et à l'\textbf{École Polytechnique de Yaoundé}, je souhaite contribuer à vos travaux de recherche.

Mon expérience la plus directement liée à votre sujet est mon \textbf{projet UROP sur l'analyse de signaux ECG}. J'y ai développé un modèle de détection d'anomalies cardiaques à partir d'images d'électrocardiogrammes : filtrage du bruit, extraction de features sur séries temporelles, et classification avec TensorFlow. Ce travail m'a appris à traiter des données réelles, souvent bruitées, et à comparer plusieurs approches de modélisation.

J'ai également participé à plusieurs compétitions de Machine Learning où j'ai dû mener des états de l'art et tester différents modèles. Pour le projet \textit{Women Safe} (détection de contenus sexistes), j'ai comparé BERT et GPT-2 sur des métriques de performance. Pour \textit{Violence Detection}, j'ai développé un modèle de classification vidéo temps réel qui m'a valu une 2\textsuperscript{ème} place. Ces expériences m'ont habitué à la démarche que vous décrivez : état de l'art, tests de plusieurs modèles, comparaison sur données réelles.

Mon passage au \textbf{Mila} comme Assistant de Recherche m'a permis de travailler avec rigueur sur des problématiques d'apprentissage profond, en reproduisant des expériences et en analysant le comportement de modèles. Ma formation à l'ENSEA en traitement du signal complète ce profil par des bases solides en filtrage et analyse spectrale.

Je suis disponible dès \textbf{avril 2026} pour un stage de 4 à 6 mois et je suis prêt à m'installer à Gramat pour la durée du stage. Je serais ravi d'échanger sur vos problématiques de détection de défauts lors d'un entretien.

Je vous prie d'agréer, Madame, Monsieur, l'expression de mes salutations distinguées.

\vspace{1cm}

\textbf{Loïc Aron MBASSI EWOLO}\\
\textit{Élève-Ingénieur ENSEA / Polytechnique Yaoundé}

\end{document}
